% 進捗報告用の全体まとめ（修論 Method / Results and Discussion の日本語版）— jresume テンプレ版
% 2026-10-12 改訂: 加藤先生の注釈 78 件を反映（対応表 docs/加藤先生コメント対応_2026-08-31報告.md）。
% 本文は docs/progress_report_2026-08-31_methods_results.tex（ltjsarticle 版）と同一（そちらから生成した）。
% 2026-09-28 訂正: 旧「文献掲載」の定義は生成コードと矛盾していたため単一出典に改めた（修論 §4.1–4.2 と同じ）。
% ビルド: このディレクトリで latexmk（latexmkrc: platex → dvipdfmx）。図は figure/*.pdf（ベクター）。
\documentclass[a4j,dvipdfmx,11pt]{jresume}

\usepackage{graphicx}
\usepackage[top=24truemm,bottom=24truemm,left=24truemm,right=24truemm]{geometry}
\usepackage{amsmath,amssymb,array}
\usepackage{booktabs}
\usepackage{enumitem}
\usepackage{xcolor}
\usepackage{tikz}
\usetikzlibrary{arrows.meta,positioning}
\usepackage[round]{natbib}
\usepackage{hyperref}

\hypersetup{colorlinks=true, linkcolor=blue!55!black, urlcolor=blue!55!black, citecolor=blue!55!black}
\setlength{\abovecaptionskip}{4pt}
\setlength{\belowcaptionskip}{0pt}
\newcommand{\rk}{\textrm{Recall@}K}
\newcommand{\Vin}{V_{\mathrm{in}}}
\newcommand{\Vout}{V_{\mathrm{out}}}
\newcommand{\Ndof}{N_{\mathrm{DoF}}}
\setlist{nosep}

% 日本語本文に合わせ、図表キャプションの見出しを日本語にする
\renewcommand{\figurename}{図}
\renewcommand{\tablename}{表}

\pagestyle{myresume}

\begin{document}
%%%%%     DATE     %%%%%
\西暦
\markright{{\footnotesize 研究進捗報告 \hfill 2026 年 8 月 31 日（10 月 12 日改訂）}}

%%%%%     TITLE     %%%%%
\medskip
\begin{center}
  {\Large 物理モデルの自動構築：手法、結果、考察の全体まとめ}
\end{center}

%%%%%     AUTHOR     %%%%%
\begin{flushright}
作成者：M2\quad 一洋
\end{flushright}
\medskip

%======================================================================
\section{本報告の位置づけ}
%======================================================================
本研究の目的は、複数の情報源から収集した数式の中から、作りたいモデルに使える可能性の高い数式を提示する手法を開発することである。提示した数式からモデルを組み立てて確かめるのは利用者である。

本報告は、問題設定と提案手法（\S\ref{sec:method}）、実験方法（\S\ref{sec:data}）、結果と考察（\S\ref{sec:overall}〜\S\ref{sec:error}）、現時点の限界（\S\ref{sec:limits}）をまとめたものである。8 月 31 日版に頂いたコメントを受けて、次の 5 点を改めた。(1) 問題設定を \citet{CombineEq} の定式化に沿って書き直し、正解と自由度を定義した。(2) 手法を全体像から説明し、各特徴量と損失関数を式で定義した。(3) 比較手法（古典的情報検索手法）の重みを 12 通り比べ、最良の重みで比較し直した。提案手法の第 1 段の重みを変えた結果も、全体比較の表に入れた。(4) 自由度ゼロ停止で選んだ式の集合は、解けるとは限らないことを明記した。(5) これまでの MAP は、並べた中にある正解式だけで平均しており、標準の定義と違っていたので、標準の定義で数え直した。

\S\ref{sec:mech}〜\S\ref{sec:topic} の分析は、第 1 段の検索の重みを文章類似度 0.7、変数の一致度 0.3 として行なったものである（重みの見直しは \S\ref{sec:redesign}）。

%======================================================================
\section{手法}\label{sec:method}
%======================================================================
\subsection{問題設定}\label{sec:problem}
作りたい物理モデルとは、値が分かっている変数（入力変数）から、求めたい変数（出力変数）を計算できる式の組である。\citet{CombineEq} は、文献から抽出した式を組み合わせてこのようなモデルを作る問題を定式化し、作りたいモデルの式の集合 $S$ が満たすべき条件として、次の 4 つの要件を定めた。ここで $V(S)$ は $S$ の式に現れる変数の集合、$\Vin$ と $\Vout$ は入力変数と出力変数の集合である。
\begin{enumerate}[leftmargin=2.2em]
  \item 入出力変数の包含：入出力変数がすべて $S$ の式に現れる（$\Vin \cup \Vout \subseteq V(S)$）。
  \item 自由度の整合：$S$ の自由度 $\Ndof(S) = |V(S)| - |S|$（変数の数と式の数の差。\citealp{Seborg}）が、入力変数の数に等しい（$\Ndof(S) = |\Vin|$）。要件 1 のもとでは、入力変数の値を与えたときに残る未知の変数の数 $|V(S) \setminus \Vin|$ が、式の数 $|S|$ に等しいことと同じである。
  \item 可解性：入出力変数以外の変数は、$S$ の 2 本以上の式に現れる。
  \item 定数の一意性：定数を表す変数は、1 本の式（変数 $=$ 定数）だけで定義される。
\end{enumerate}
これらは必要条件であり、満たした式の組がモデルとして役に立つとは限らない~\citep{CombineEq}。

\citet{CombineEq} は、与えられた式の集合から、要件を満たす組合せをすべて求める。本研究はこれと異なり、多数の文献から作った大規模な数式データベースから、作りたいモデルに使える可能性の高い式を選んで提示する。入力はクエリ（検索の入力として与える文や語）$q = (c, \Vin, \Vout)$ であり、$c$ は作りたいモデルを自然言語で説明した文（説明文）、$\Vin$ と $\Vout$ は入力変数と出力変数の記号の集合である。$V_q = \Vin \cup \Vout$ を入出力変数とよぶ。数式データベース $E = \{e_1, \dots, e_N\}$ の各式 $e$ は、次の 5 つの情報を持つ（作り方は \S\ref{sec:db}）。
\begin{itemize}[leftmargin=1.5em]
  \item 式の表記：\LaTeX{} 形式の式。
  \item 変数の集合 $V(e)$：式に現れる変数記号の集合。各記号には、文献に書かれた定義の文章が付く。
  \item 式周辺の文章：式を抽出するときに LLM が、式の前後の本文をもとに書き出した、その式の物理的な意味を説明する文章。
  \item 出典：式を抽出した文献の名前。
  \item 分野ラベル $d(e)$：式が扱う分野やプロセスの名前（例：CSTR）。式を抽出するときに LLM が自由に書いた文字列であり、決まった分類から選んだものではない。
\end{itemize}
変数どうしは、記号の文字列が完全に一致するときだけ同じ変数とみなす。この照合は、クエリの変数と式の変数を比べるときも、式どうしの変数を比べるときも同じであり、表記の統一はしない（$C_A$ と $C_A(t)$ は別の変数として扱う）。

手法の出力は、データベースの式の順位である。提案手法は、作りたいモデルに使える可能性の高い式ほど上位に並べ、上位の式を提示する。何本を提示するかは \S\ref{sec:dofstop} で述べる。

手法を評価するため、評価に使うケースには、クエリとともに、作りたいモデルを構成する式の集合 $R_q \subset E$ を 1 つ定めてある（\textbf{正解集合}。作り方は \S\ref{sec:cases}）。$R_q$ に含まれる式を\textbf{正解式}、それ以外の式を\textbf{不正解式}とよび、$K = |R_q|$ を正解式の数とする。評価に使う 1{,}823 ケースの正解集合は、すべて要件 1〜3 を満たす（要件 4 は、どの変数が定数を表すかをデータベースが区別していないため、確かめていない）。ただし要件 1〜3 は変数と式の数え方の条件であり、未知の変数のそれぞれに、それを決める式を 1 本ずつ割り当てられるかは見ない。出典横断の 91 件（1{,}823 件の 5.0\%）は、この割り当てができず、入力変数を与えても一部の未知の変数が決まらない（どれかの正解式の変数がすべて入力変数になっている。91 件のうち 87 件は、入出力変数を選び直した派生のケース）。式は ID（出典と出典内の番号の組）で区別するので、正解式と同じ式が別の出典にあり、それを選んだ場合は、検索手法の評価では不正解と数える（\S\ref{sec:eval}、\S\ref{sec:errtypes}）。

\subsection{提案手法の概要}\label{sec:overview}
提案手法は、2 段階の検索でデータベースの式を並べ替える（図~\ref{fig:overview}）。
\begin{itemize}[leftmargin=1.5em]
  \item \textbf{第 1 段の検索（候補の絞り込み）}：データベースの全 $N$ 式に、クエリとの関連の強さを表す簡単な点数を付け、上位 $N_1$ 件を候補として残す。
  \item \textbf{第 2 段の検索（候補の並べ直し）}：候補 $N_1$ 件のそれぞれについて 10 個の特徴量を計算し、多層パーセプトロン（multilayer perceptron, MLP）が特徴量から候補の点数を計算する。候補を点数の高い順に並べ直す。10 個のうち 3 個は、候補をほかの式の集合（\textbf{参照集合}）と比べる特徴量である。物理モデルでは、ある式に現れる変数をほかの式が決めるので、ある式が使えるかどうかは、一緒に提示する式によって変わる。この 3 個はそれを点数に反映させる。
\end{itemize}
第 2 段の並べ直し方には 2 通りある。基本形の\textbf{一括の並べ直し}は、参照集合を第 1 段の上位 $m$ 件に固定し、全候補を 1 回で点数化する（\S\ref{sec:stage2}）。拡張形の\textbf{逐次選択}は、式を 1 本ずつ選び、選ぶたびに、それまでに選んだ式を参照集合にして点数を計算し直す（\S\ref{sec:seqmethod}）。正解式の数 $K$ は事前に分からないので、何本を提示するかは、逐次選択に停止規則を加えて決める（\S\ref{sec:dofstop}）。手法の定数（$N_1$、$m$、重みなど）の値は \S\ref{sec:compared} に示す。

\begin{figure}[t]\centering
\begin{tikzpicture}[
  box/.style={draw, rounded corners=2pt, align=center, font=\footnotesize, inner sep=4pt, minimum height=24mm},
  arr/.style={-{Stealth[length=2.4mm]}, thick}]
  \node[box, text width=27mm] (in) {\textbf{クエリ} $q$\\説明文 $c$\\入力変数 $\Vin$\\出力変数 $\Vout$\\[3pt]\textbf{データベース} $E$\\（$N$ 式）};
  \node[box, text width=33mm, right=6mm of in] (s1) {\textbf{第 1 段の検索}\\全 $N$ 式を点数 $s_1$ で\\並べ、上位 $N_1$ 件を\\候補として残す\\[3pt]$s_1$：文章類似度と\\変数の一致度の重み付き和};
  \node[box, text width=40mm, right=6mm of s1] (s2) {\textbf{第 2 段の検索}\\候補ごとに特徴量 10 個\\（うち集合の特徴量 3 個）\\を計算し、MLP の点数で\\候補を並べ直す\\[3pt]一括の並べ直し、\\または逐次選択};
  \node[box, text width=31mm, right=6mm of s2] (out) {\textbf{提示}\\評価：上位 $K$ 件\\[3pt]利用時：\\自由度ゼロ停止までに\\選んだ式};
  \draw[arr] (in) -- (s1);
  \draw[arr] (s1) -- (s2);
  \draw[arr] (s2) -- (out);
\end{tikzpicture}
\caption{提案手法の全体像。第 1 段の検索で候補を $N_1$ 件に絞り、第 2 段の検索で候補を並べ直して上位の式を提示する。}
\label{fig:overview}
\end{figure}

\subsubsection{第 1 段の検索：候補の絞り込み}\label{sec:stage1}
第 1 段の検索は、データベースの各式 $e$ に対して、クエリ $q$ との点数
\begin{equation}
  s_1(q, e) = w_{\mathrm{t}} \cos\bigl(\phi(c), \phi(T(e))\bigr) + w_{\mathrm{v}} \, J\bigl(V_q, V(e)\bigr)
  \label{eq:stage1}
\end{equation}
を計算し、点数の高い上位 $N_1$ 件を候補 $C_q$ とする。第 1 項は説明文と式の文章の近さ（\textbf{文章類似度}）、第 2 項は入出力変数と式の変数の重なり（\textbf{変数の一致度}）であり、$w_{\mathrm{t}}$ と $w_{\mathrm{v}}$ はその重みである。$T(e)$ は、式 $e$ の式周辺の文章、式の表記、分野ラベルをつなげた文字列である。$\phi(\cdot)$ は文字列の TF-IDF（term frequency--inverse document frequency）ベクトルで、各単語に、その文字列での出現回数と、データベースの式の文字列全体での珍しさを掛けた重みを付ける。単語 1 語と、隣り合う 2 語の組を単位とし、語彙はデータベースの全式の $T(e)$ から作る。$\cos(\cdot, \cdot)$ は 2 つのベクトルのコサイン類似度、$J(A, B) = |A \cap B| / |A \cup B|$ は 2 つの集合の Jaccard 係数である。

第 2 段の検索は候補の中でしか並べ直せないので、正解式のうち候補に入った割合（\textbf{被覆率}、\S\ref{sec:eval}）が、第 2 段で得られる精度の上限になる。被覆率は、重み $(w_{\mathrm{t}}, w_{\mathrm{v}})$ と候補の件数 $N_1$ の両方で変わる（\S\ref{sec:redesign}）。

\subsubsection{第 2 段の検索：一括の並べ直し}\label{sec:stage2}
第 2 段の検索は、候補 $e \in C_q$ ごとに 10 個の特徴量を並べたベクトル $x(q, e)$ を作り、MLP $f$ で点数 $f(x(q, e))$ を計算して、候補を点数の高い順に並べ直す。表~\ref{tab:features} に 10 個の特徴量をまとめる。以下で各特徴量を定義する。

\textbf{基本の 7 個}は、クエリと候補 1 本を比べる。
\begin{itemize}[leftmargin=1.5em]
  \item $x_1$ 文章類似度：式~\eqref{eq:stage1} の第 1 項と同じ $\cos(\phi(c), \phi(T(e)))$。
  \item $x_2$ 変数の一致度：式~\eqref{eq:stage1} の第 2 項と同じ $J(V_q, V(e))$。
  \item $x_3$ 潜在意味類似度：データベースの全式の $T(e)$ の TF-IDF ベクトルに切断特異値分解（truncated singular value decomposition, SVD）を当てはめて 256 次元に縮約し（潜在意味解析（latent semantic analysis, LSA））、縮約した空間で、クエリと式 $e$ のコサイン類似度をとる。クエリ側の文字列には、説明文 $c$ に入出力変数の記号を付け加える。縮約した空間では、同じ文字列に一緒に現れやすい単語が同じ次元にまとまるので、同じ単語を含まなくても、似た文字列に現れる単語を含めば値が大きくなる。
  \item $x_4$ 入力変数の被覆率：$|\Vin \cap V(e)| / |\Vin|$。入力変数のうち式に現れる割合。
  \item $x_5$ 出力変数の被覆率：$|\Vout \cap V(e)| / |\Vout|$。出力変数のうち式に現れる割合。
  \item $x_6$ 式の特化度：$|V_q \cap V(e)| / |V(e)|$。式の変数のうち入出力変数が占める割合。入出力変数以外の変数を持たない式ほど 1 に近い。
  \item $x_7$ 分野ラベルの一致：分野ラベル $d(e)$ が説明文 $c$ に文字列として含まれるか、$c$ が $d(e)$ に含まれれば 1、そうでなければ 0（大文字と小文字は区別しない）。
\end{itemize}

\textbf{集合の 3 個}は、候補 $e$ を参照集合 $S$ と比べる。参照集合は、候補を評価するときに比べる相手とする式の集合であり、一括の並べ直しでは、第 1 段の検索の結果得られる上位 $m$ 件とする。候補自身が参照集合に入っているときは、その候補を除いた参照集合と比べる。$V(S)$ は $S$ の式に現れる変数の集合である。
\begin{itemize}[leftmargin=1.5em]
  \item $x_8$ 補完性：$|(V_q \cap V(e)) \setminus V(S)| / |V_q|$。入出力変数のうち、参照集合の式が持つ変数の集合には含まれないが、候補式が持つ変数の集合には含まれる変数の割合。
  \item $x_9$ 一貫性：$|V(e) \cap V(S)| / |V(e)|$。候補式の変数のうち、参照集合の式にも現れる変数の割合。
  \item $x_{10}$ 分野ラベルの一致度（集合）：$|\{e' \in S : d(e') = d(e)\}| / |S|$。参照集合の式のうち、分野ラベルが候補式の分野ラベルと文字列として完全に一致する式の割合。
\end{itemize}

\begin{table}[t]\centering
\footnotesize
\caption{第 2 段の検索が使う 10 個の特徴量のまとめ（定義は本文）。$V_q = \Vin \cup \Vout$ は入出力変数、$V(e)$ は候補式の変数の集合、$S$ は参照集合、$T(e)$ は式の文字列、$d(e)$ は分野ラベル。「群」は、変数の重複率を測るか、話題の近さを測るかを表す。}
\label{tab:features}
\begin{tabular}{@{}llll@{}}
\toprule
 & 特徴量 & 定義 & 群 \\
\midrule
$x_1$ & 文章類似度 & $\cos(\phi(c), \phi(T(e)))$ & 話題 \\
$x_2$ & 変数の一致度 & $J(V_q, V(e))$ & 変数 \\
$x_3$ & 潜在意味類似度 & LSA で 256 次元にした空間でのコサイン類似度 & 話題 \\
$x_4$ & 入力変数の被覆率 & $|\Vin \cap V(e)| / |\Vin|$ & 変数 \\
$x_5$ & 出力変数の被覆率 & $|\Vout \cap V(e)| / |\Vout|$ & 変数 \\
$x_6$ & 式の特化度 & $|V_q \cap V(e)| / |V(e)|$ & 変数 \\
$x_7$ & 分野ラベルの一致 & $d(e)$ と $c$ の一方が他方に含まれれば 1 & 話題 \\
\midrule
$x_8$ & 補完性 & $|(V_q \cap V(e)) \setminus V(S)| / |V_q|$ & 変数 \\
$x_9$ & 一貫性 & $|V(e) \cap V(S)| / |V(e)|$ & 変数 \\
$x_{10}$ & 分野ラベルの一致度（集合） & $|\{e' \in S : d(e') = d(e)\}| / |S|$ & 話題 \\
\bottomrule
\end{tabular}
\end{table}

MLP は、10 個の特徴量を入力とし、64 ユニットの隠れ層 2 層（活性化関数 ReLU、ドロップアウト率 0.1）を経て 1 つの点数を出力する。訓練では、正解式の点数が、同じケースの不正解の候補の点数よりマージン $\gamma$ 以上大きくなるように学習する。訓練ケース $q$ の正解式 $e^+ \in P_q = R_q \cap C_q$ と不正解の候補 $e^-$ の組ごとに
\begin{equation}
  \ell(q, e^+, e^-) = \max\bigl(0,\ \gamma - f(x(q, e^+)) + f(x(q, e^-))\bigr)
  \label{eq:loss}
\end{equation}
を計算する。$e^-$ は、不正解の候補 $C_q \setminus R_q$ から無作為に選んだ $n_{\mathrm{neg}}$ 件である。訓練ケース 16 件ごとに、そこに含まれるすべての組の $\ell$ の平均を損失として、MLP の重みを更新する。正解式の点数が不正解の候補の点数を $\gamma$ 以上上回れば、その組の損失は 0 になる。

\subsubsection{第 2 段の拡張：逐次選択}\label{sec:seqmethod}
一括の並べ直しは、参照集合を第 1 段の上位 $m$ 件に固定する。正解式の数 $K$ が大きいケースでは、実際に上位に並ぶ式と、固定した参照集合が食い違っていく。逐次選択は、参照集合を「それまでに選んだ式」に置き換えて、式を 1 本ずつ選ぶ。推論の手順は次のとおりである。
\begin{enumerate}[leftmargin=2.2em]
  \item 1 本目には、一括の並べ直しの点数（固定の参照集合で計算した点数）が最も高い候補を選ぶ。
  \item 2 本目以降は、それまでに選んだ式の集合を参照集合として、残りの候補の $x_8$〜$x_{10}$ を計算し直し、点数が最も高い候補を選ぶ。
  \item 手順 2 を、20 本選ぶか候補がなくなるまで繰り返す。選ばれなかった候補は、一括の並べ直しの点数の順に後ろへ並べる。
\end{enumerate}
逐次選択には、MLP の訓練方法の異なる 2 通りがある。\textbf{逐次選択（推論のみ）}は、一括の並べ直しと同じ方法で訓練した MLP を、上の手順で使う。\textbf{逐次選択（訓練も逐次）}は、訓練でも推論と同じ形の参照集合を使う。各訓練ケースで、候補に入った正解式を第 1 段の検索の順位の順に $e^+_1, e^+_2, \dots$ と並べ、段階 $t = 1, 2, \dots$ ごとに損失を加える。段階 $t$ では、正解式 $e^+_1, \dots, e^+_{t-1}$ を「それまでに選んだ式」とみなして参照集合とし（$t = 1$ では第 1 段の上位 $m$ 件）、その参照集合で計算した特徴量について、次の正解式 $e^+_t$ と不正解の候補の組の損失 $\ell$（式~\eqref{eq:loss}）を加える。段階の数は 1 ケースあたり最大 8 である。こうすると、訓練でも推論と同じ問い（「これらの式を選んだ今、次に選ぶべき式はどれか」）を学ぶ。ただし、推論では自分が選んだ式（誤りを含みうる）を参照集合にするのに対し、訓練では正解式を参照集合にする点が異なる（\S\ref{sec:limits}）。

\subsubsection{提示する式の数の決め方：自由度ゼロ停止}\label{sec:dofstop}
評価では正解式の数 $K$ を与えて上位 $K$ 件を採るが（\S\ref{sec:eval}）、実際の利用者は $K$ を事前に知らない。そこで逐次選択で 1 本選ぶごとに、それまでに選んだ式の集合 $S$ が次の 2 つの条件を満たすかを調べ、初めて満たした時点で選択を止めて、$S$ を提示する。20 本選んでも満たさなければ、選んだ 20 本を提示する。以下、この規則を\textbf{自由度ゼロ停止}とよぶ。
\begin{enumerate}[label=(\alph*), leftmargin=2.2em]
  \item 出力変数がすべて $S$ の式に現れる（$\Vout \subseteq V(S)$）。
  \item 入力変数以外の変数（未知の変数）の数が、式の数以下である（$|V(S) \setminus \Vin| \le |S|$）。
\end{enumerate}
条件 (a) は要件 1 のうち出力変数の部分であり、条件 (b) は要件 2 の等号を不等号に緩めたものである。要件 3 と 4 は調べず、式が互いに独立か、矛盾しないかも調べない。そのため、止まった時点の $S$ が解けるとは限らない。例えば、$\Vin = \{x, z\}$、$\Vout = \{y_1, y_2\}$ のクエリで、$y_1 = x$、$y_1 = x + 1$、$y_2 = z$ の順に選ぶと、3 本目で (a) と (b) を満たして止まる（未知の変数は $y_1$ と $y_2$ の 2 個で、式は 3 本）。しかし、$y_1$ を決める 2 本の式は矛盾しており、自由度は $\Ndof(S) = 4 - 3 = 1$ で入力変数の数 2 に等しくないので、要件 2 も満たさない。自由度ゼロ停止は、提示する式の数を決めるための規則であり、提示した式でモデルが組めるかを確かめるのは利用者である。

%======================================================================
\section{実験方法}\label{sec:data}
%======================================================================
\subsection{数式データベース}\label{sec:db}
データベースは、361 の出典から集めた $11{,}146$ 式である。出典のうち 329 件は文献で、$10{,}476$ 式を含む。残る 32 件は、LLM（Gemini 2.5 Pro）が分野ごとに教科書の水準の基本式を生成した式集（\textbf{基礎式集}、670 式）であり、元になる文献はない。

文献は 2 つの方法で集めた。研究の初期には、化学工学の教科書と論文を手作業で集めた。その後、データベースを化学工学以外にも広げるため、反応工学、伝熱、電気化学、バイオプロセス、音響学など 32 分野のそれぞれについて検索語を決め、学術文献データベース OpenAlex から、オープンアクセスで PDF を入手できる文献（2006 年以降に出版され、被引用数が 10 回を超える論文、総説、書籍の章など）を、検索語との関連度の高い順に集めた。1 分野あたり 10 報を目標とし、足りない分野は被引用数の条件を緩めて補った。

% 抽出モデル: 初期の 39 報は Gemini 2.5 Pro（extract_equations.py）、追加した文献は Claude（extract_equations_claude.py、
%   本人確認 2026-10-12「追加した文献の抽出はClaudeを使った」）。Claude の版は記録がない（スクリプトの既定は claude-opus-4-5-20250805、docstring は Opus 4.7）。
式の抽出には LLM を使い、研究の初期に集めた 39 報は Gemini 2.5 Pro で、その後に追加した文献は Claude で抽出した。どちらも各文献の PDF を読ませ、番号の付いた式（本文と付録の式）を \LaTeX{} 形式で書き出させ、あわせて式ごとに、変数記号とその定義、式の前後の本文をもとに式の物理的な意味を説明する文章（式周辺の文章）、分野ラベルを、決めた形式（JSON）で返させた。Claude には、数値の計算例や演習問題の式は除くよう指示した。変数記号は文献の表記のまま書き出させ、表記の統一はしていない。その結果、データベースには異なる変数記号が $13{,}220$ 種、異なる分野ラベルが $1{,}998$ 種ある。

\subsection{評価ケース}\label{sec:cases}
評価ケースは、クエリ（説明文と入出力変数）と正解集合の組である。正解集合の作り方が異なる 3 種類、計 1{,}823 件を作った（表~\ref{tab:cases}）。どの種類でも、入出力変数は正解式に現れる変数記号から選んでいる。

\textbf{単一出典}（92 件）：LLM（Gemini 2.5 Pro）に、データベースの式の一覧を渡してケースを作らせた。一覧には、式ごとに ID（出典名を含む）、分野ラベル、変数記号だけを載せ、式の表記と式周辺の文章は載せていない。LLM には、一覧の式を組み合わせて、物理的に意味のあるケース（説明文、入力変数、出力変数、正解集合）を作るよう指示し、正解集合の条件として、入力変数の値から出力変数をすべて求められることと、必要十分かつ最小であること（余分な式を含まないこと）を課した。LLM が作った 123 件のうち、正解式がすべてデータベースにあり、入出力変数がすべて正解式に現れ、入力変数を与えたときの未知の変数の数と式の数が等しい 92 件を残した。したがって、92 件という数は出典の数（361 件）とは対応しない。92 件のうち 91 件の正解集合は 1 つの出典の式だけから成るので、単一出典とよぶ。正解集合は、出典が 1 つのモデルとして示す式の組ではなく、LLM が最小と判断した組であり、92 件中 67 件は 1 式、61 件は基礎式集の式だけから成る。

\textbf{出典横断}（731 件）：異なる出典の式を 2〜3 本組み合わせて作った。500 件はプログラムで作った。分野ラベルを反応器や伝熱などの 10 の分類にまとめ、組み合わせてよい分類の組をあらかじめ決めておき、その組に属する異なる出典の 2 式で、変数記号を 2 個以上共有し、共有する記号の定義の文章が食い違わないものを列挙した（微分方程式を含む組を優先した）。1 組の 2 式から、それぞれの式だけに現れる変数を 1 個ずつ選んで出力変数とし、残りの変数を入力変数とした（1 組から最大 5 通り）。説明文は、分類ごとに決めた文型に当てはめて作った。残る 231 件は、開発の途中で LLM（Claude）が書いた 33 組と、その入出力変数を選び直した 198 件であり、33 組は、入力変数を与えたときの未知の変数の数と式の数が等しいことを確かめている。この 231 件のうち 189 件は、式どうしが変数記号を 1 つも共有しない。

\textbf{合成}（1{,}000 件）：正解集合の大きさ $X$（正解集合に含める式の本数の目標値。1〜10）を決め、規則で正解集合を作った。まず、ある出典から微分方程式を 1 本選び、微分される変数を出力変数とする。次に、それまでの式に現れる出力変数以外の変数を含む式を、同じ出典を優先して 1 本加え、その式の変数から新しい出力変数を 1 個選ぶ。これを式が $X$ 本になるまで繰り返し、出力変数に選ばれなかった変数をすべて入力変数とする。こうして作った正解集合では、式の数と出力変数の数が等しく、入力変数を与えると未知の変数の数と式の数が等しい。$X = 1$〜10 のそれぞれで 100 件作り、説明文だけを LLM（Claude Sonnet 4.5）に書かせた。

\begin{table}[t]\centering
\small
\caption{評価ケースの 3 種類（作り方は本文）。「式数」は正解集合の式の本数の平均（括弧は最大）、「出典数」は正解集合が含む出典の数の平均。}
\label{tab:cases}
\begin{tabular}{@{}lp{0.52\linewidth}rrr@{}}
\toprule
種類 & 正解集合の作り方 & 件数 & 式数 & 出典数 \\
\midrule
単一出典 & LLM が、入出力を決めるのに必要十分かつ最小と判断した組 & 92 & 1.34 (4) & 1.01 \\
出典横断 & 異なる出典の式を 2〜3 本組み合わせた組 & 731 & 2.03 (3) & 2.02 \\
合成 & 微分方程式から始め、規則で $X$ 本（$X = 1$〜10）に育てた組 & 1{,}000 & 5.50 (10) & 1.10 \\
\bottomrule
\end{tabular}
\end{table}

このほか、単一出典のケースから、既存ケースの入出力の入替や説明文の言い換えによって、拡張ケース 1{,}015 件を作成した。738 件は、正解集合はそのままで入出力変数を選び直したもの（正解式の変数から、式の本数と同じ数の出力変数を無作為に選ぶ）、276 件は説明文を決まった言い換え規則で書き換えたもの（LLM は使っていない）、1 件は入力変数と出力変数を入れ替えたものである。拡張ケースは訓練だけに使い、評価からは除いた。拡張ケースは元のケースとほとんど同じ易しいケースなので、評価に入れると、古典的情報検索手法（重み $(0.7, 0.3)$、\S\ref{sec:compared}）の $\rk$ が 0.521 から 0.622 に上がり、提案手法との差が $+0.222$ から $+0.156$ に縮んで見える。

評価は、3 種類すべてを使う\textbf{設定 A}（1{,}823 件）と、合成だけを使う\textbf{設定 B}（1{,}000 件）で行なう。設定 B は、正解式の数が多いケースだけで評価するための設定である（正解式の数の平均は、設定 A で 3.90、設定 B で 5.50）。設定 A では 90.8\% のケースで正解式が 2 本以上ある。

\subsection{訓練とテストの分割}\label{sec:split}
訓練、検証、テストへの分割は、出典を単位として行なう。同じ出典の式を正解に含むケースが訓練とテストの両方に入ると、テストの正解を訓練で見たことになるためである。各ケースは、正解式の出典のうち最も多いものに属するとする。出典の一覧を無作為に並べ替え、出典の 20\% をテスト、20\% を検証、60\% を訓練に割り当てる。検証用の出典は、今回の実験では使っていない。

この割り当てを 1 回だけ無作為に行なうと、正解式の数の多いケースが一方に偏ることがある。そこで、割り当てを 400 通り無作為に作り、その中から、訓練とテストでケースの分布が最もそろう 1 通りを選ぶ（\textbf{層化}）。分布をそろえる量は、ケースごとの正解式の数、入力変数の数、出力変数の数、正解式の出典の数の 4 つである。それぞれの量について、訓練とテストの度数分布を割合に直し、値ごとの割合の差の絶対値を足した L1 距離を求める。4 つの量の L1 距離の和が最小になる割り当てを選ぶ。

この分割を、乱数の種を 10 通り（42、123、456、789、1024、2024、3141、5926、7777、9999）に変えて繰り返す。乱数の種が変わると、400 通りの候補も、選ばれる分割も変わるので、10 通りの異なる分割で評価することになる。MLP の初期値と、不正解の候補の抽出も、同じ乱数の種で決まる。テストのケースは分割ごとに 394〜559 件である（設定 A）。すべての手法を同じ 10 通りの分割で評価し、10 通りの平均と標準偏差を報告する。2 つの手法の差は、10 通りの分割の値を対にした Wilcoxon の符号順位検定で調べる。10 対の検定で $p$ 値の最小は 0.00195 であり、これは 10 通りすべてで一方の手法が上回ったことを意味する。

\subsection{比較手法と手法の定数}\label{sec:compared}
比較手法は次の 2 つである。

\textbf{古典的情報検索手法（baseline）}：第 1 段の検索と同じ点数 $s_1$（式~\eqref{eq:stage1}）で、データベースの全式を並べる。学習も候補の絞り込みもしない。重み $(w_{\mathrm{t}}, w_{\mathrm{v}})$ は、$(1.0, 0.0)$ から $(0.0, 1.0)$ までの 12 通りを比べ（\ref{app:weights}）、テストの $\rk$（\S\ref{sec:eval}）の 10 通りの平均が最も高い $(0.4, 0.6)$ を本文で用いる（設定 A と B で同じ）。テストの結果で重みを選んでいるので、baseline に有利な比較である。文章類似度を BM25~\citep{BM25} や E5（文の埋め込みベクトル）~\citep{E5} の類似度に置き換えた比較も\ref{app:weights} に示す。

\textbf{LLM 直接生成}：LLM（Claude Sonnet 4.5）に、説明文と入出力変数を与え、出力変数を求めるのに必要な式を、重要な順に 1 行に 1 本ずつ \LaTeX{} 形式で書かせる（最大 $\max(12, K + 4)$ 本）。判定（\S\ref{sec:eval}）に費用がかかるため、各設定から無作為に選んだ 50 件を 1 回だけ評価した。

\textbf{提案手法}：\S\ref{sec:overview} の手法である。候補の件数 $N_1 = 50$、参照集合の大きさ $m = 5$、マージン $\gamma = 0.1$、不正解の抽出数 $n_{\mathrm{neg}} = 8$ とし、AdamW（学習率 $10^{-3}$、重み減衰 $10^{-4}$）で 15 エポック訓練する。第 1 段の重みは、以前から使っていた $(0.7, 0.3)$ と、4 通り（$(0.7, 0.3)$、$(0.5, 0.5)$、$(0.3, 0.7)$、$(0.0, 1.0)$）を比べて最良だった $(0.3, 0.7)$（\S\ref{sec:redesign}）の 2 つの結果を示す。逐次選択と自由度ゼロ停止の実験では、候補の件数を $N_1 = 200$ に広げる（理由は \S\ref{sec:greedy}）。

\subsection{評価方法}\label{sec:eval}
\textbf{正解の判定}：検索手法（baseline と提案手法）はデータベースの式を返すので、返した式の ID が正解集合 $R_q$ に含まれるかで、正解かどうかを判定する。LLM 直接生成はデータベースにない文字列の式を書くので、ID では判定できない。そこで、書かれた式（先頭 12 本まで）のそれぞれが、どの正解式と数学的に等価かを、生成とは別の LLM（Claude Opus 4.8）に判定させ、等価と判定された式を正解と数える。判定の結果が得られなかったケースは除いた（設定 A で 50 件中 2 件）。この判定を確かめるため、別の LLM（Claude Fable 5.1）に同じ指示で、3 つの標本（設定 A、設定 B、拡張を含む全ケース）の 608 組をすべて判定し直させた。2 つの判定は 88.0\% の組で一致し（Cohen の $\kappa = 0.72$）、食い違った 73 組のうち 71 組は、Claude Opus 4.8 が等価とし、Claude Fable 5.1 が等価でないとした組であった。人手の判定との照合と、同じ判定を繰り返したときのばらつきは、まだ調べていない（\S\ref{sec:limits}）。

ID による判定では、正解式と同じ式が別の出典にあり、それを選んだ場合も不正解と数える。等価かどうかで判定する LLM 直接生成と比べると、検索手法に厳しい判定である。正解式 3{,}132 本のうち 386 本（12.3\%）は、空白などを除いて同じ文字列の式が別の出典にもある（\S\ref{sec:errtypes}）。

\textbf{指標}：各ケースについて次の量を求め、テストのケースで平均した値を 1 つの分割の値とする。$\mathrm{Top}_n(q)$ は、手法が並べた上位 $n$ 件である。
\begin{itemize}[leftmargin=1.5em]
  \item $\rk = |R_q \cap \mathrm{Top}_K(q)| / K$（主指標）：上位 $K$ 件（$K$ は正解式の数）に入った正解式の割合。上位 $K$ 件のうちの正解式の割合とも等しい（情報検索でいう R-precision）。$K$ がケースごとに 1〜10 と異なるので、見る件数を $K$ に合わせる。
  \item MAP（mean average precision）：各ケースの平均適合率 $\mathrm{AP}(q) = \frac{1}{K} \sum_{e \in R_q} P(e)$ のケース平均。$P(e)$ は、正解式 $e$ の順位までの上位の式のうち正解式の割合である（$e$ が並べた中にないときは 0。提案手法では、第 1 段の候補に入らなかった正解式がこれにあたる）。正解式がどれだけ上位に集まっているかを表す。
  \item Recall@20 $= |R_q \cap \mathrm{Top}_{20}(q)| / K$：利用者が上位 20 件を見る場合に、正解式がどれだけ含まれるか。
  \item 被覆率 $= |R_q \cap C_q| / K$：第 1 段の検索の候補に入った正解式の割合。
\end{itemize}
LLM 直接生成の $\rk$ は、正解式のうち、書かれた先頭 $K$ 本の中に等価な式がある割合とする。あわせて、順位を問わない被覆率（書かれた式全体の中に等価な式がある正解式の割合）も示す。LLM 直接生成の MAP と Recall@20 は求めていない。

自由度ゼロ停止は、式の順位ではなく、式の集合 $S$ を選ぶ。そこで、完全一致率（$S = R_q$ となったケースの割合）、集合 F1（$2|S \cap R_q| / (|S| + |R_q|)$ のケース平均）、停止条件の充足率（$S$ が \S\ref{sec:dofstop} の条件 (a) と (b) を満たしたケースの割合）で評価し、同じ逐次選択で上位 $K$ 本を採った場合と比べる。

%======================================================================
\section{結果 1：全体比較 --- 正解式と入力変数が多いケースほど差が大きい}\label{sec:overall}
%======================================================================
表~\ref{tab:overall} に全体比較を示す。最良の重みの baseline に比べ、提案手法（第 1 段の重み $(0.7, 0.3)$）は $\rk$ を設定 A で 0.607 から 0.743 に（$+0.136$）、設定 B で 0.559 から 0.689 に（$+0.130$）上げ、どちらの設定でも 10 通りの分割すべてで上回った（Wilcoxon の符号順位検定、$p = 0.00195$）。第 1 段の重みを $(0.3, 0.7)$ にすると、$\rk$ は 0.780（設定 A）と 0.740（設定 B）になり、最良の重みの baseline を $+0.173$ と $+0.181$ 上回る。この重みの提案手法は、$\rk$、MAP、Recall@20 のいずれでも、どちらの設定でも、10 通りの分割すべてで最良の重みの baseline を上回った（$p = 0.00195$）。

一方、重み $(0.7, 0.3)$ の提案手法の Recall@20 は、最良の重みの baseline とほとんど変わらない（設定 A で 0.909 と 0.907、$p = 0.49$。設定 B で 0.838 と 0.827、$p = 0.049$）。提案手法は第 1 段の候補 50 件の中でしか並べ直せず、重み $(0.7, 0.3)$ の候補には正解式の 92.0\% しか入らないためである（設定 A）。第 1 段の重みを $(0.3, 0.7)$ にすると、候補に入る正解式が 96.7\% に増え（設定 A）、Recall@20 は 0.957（設定 A）と 0.925（設定 B）に上がる。

\begin{table}[t]\centering
\caption{全体比較（検索手法は 10 通りの分割の平均 $\pm$ 標準偏差）。「重み」は第 1 段の点数 $s_1$ の重み（文章類似度、変数の一致度）。baseline の $(0.4, 0.6)$ は 12 通りの重みで最良のもの（\ref{app:weights}）。MAP は標準の定義（\S\ref{sec:eval}）による。8 月 31 日版の MAP（提案手法で 0.863 と 0.832）は、並べた中にある正解式だけで平均した値で、第 1 段の候補に入らなかった正解式が除かれていた。LLM 直接生成は各設定 50 件の無作為標本の 1 回評価で、括弧内は判定器を Claude Fable 5.1 に替えた値。}
\label{tab:overall}
\footnotesize
\setlength{\tabcolsep}{4pt}
\begin{tabular}{@{}lllccc@{}}
\toprule
設定 & 手法 & 重み & $\rk$ & MAP & Recall@20 \\
\midrule
A & 古典的情報検索手法（baseline） & $(0.4, 0.6)$ & $0.607 \pm 0.043$ & $0.692 \pm 0.046$ & $0.907 \pm 0.016$ \\
（1{,}823 件） & 古典的情報検索手法（baseline） & $(0.7, 0.3)$ & $0.521 \pm 0.045$ & $0.581 \pm 0.045$ & $0.836 \pm 0.023$ \\
 & LLM 直接生成（$n = 48$） & --- & 0.254 (0.198) & --- & --- \\
 & 提案手法（一括の並べ直し） & $(0.7, 0.3)$ & $0.743 \pm 0.054$ & $0.796 \pm 0.043$ & $0.909 \pm 0.014$ \\
 & 提案手法（一括の並べ直し） & $(0.3, 0.7)$ & $\mathbf{0.780 \pm 0.048}$ & $\mathbf{0.850 \pm 0.038}$ & $\mathbf{0.957 \pm 0.011}$ \\
\midrule
B & 古典的情報検索手法（baseline） & $(0.4, 0.6)$ & $0.559 \pm 0.038$ & $0.627 \pm 0.036$ & $0.827 \pm 0.027$ \\
（合成 1{,}000 件） & 古典的情報検索手法（baseline） & $(0.7, 0.3)$ & $0.466 \pm 0.037$ & $0.512 \pm 0.035$ & $0.721 \pm 0.032$ \\
 & LLM 直接生成（$n = 50$） & --- & 0.226 (0.150) & --- & --- \\
 & 提案手法（一括の並べ直し） & $(0.7, 0.3)$ & $0.689 \pm 0.033$ & $0.715 \pm 0.034$ & $0.838 \pm 0.026$ \\
 & 提案手法（一括の並べ直し） & $(0.3, 0.7)$ & $\mathbf{0.740 \pm 0.040}$ & $\mathbf{0.795 \pm 0.040}$ & $\mathbf{0.925 \pm 0.020}$ \\
\bottomrule
\end{tabular}
\end{table}

LLM 直接生成の $\rk$ は 0.254（設定 A）と 0.226（設定 B）である。LLM は 1 ケースあたり平均 11.6 本の式を書いた（正解式の数は、標本の平均で 3.9 本と 5.6 本）。書かれた式全体では、正解式の 41.9\%（設定 A）と 34.3\%（設定 B）に等価な式があった（順位を問わない被覆率）が、そのうち重要な順に並べた先頭 $K$ 本に入っていたのは一部である（$\rk$ 0.254 と 0.226）。等価の判定器を Claude Fable 5.1 に替えると、$\rk$ は 0.198 と 0.150、順位を問わない被覆率は 0.338 と 0.234 に下がる。この比較は 50 件程度を 1 回だけ評価したものであり、全ケースでの評価と複数回の評価、文献を与えた LLM との比較は今後の課題である（\S\ref{sec:summary}）。

表~\ref{tab:overall} の差を、正解式の数と入力変数の数で分けて見る（設定 A、10 通りの分割のテストのケースを合わせた延べ 4{,}594 件。\ref{app:bucket}）。正解式が 1 本のケースでは、提案手法（重み $(0.7, 0.3)$）の $\rk$ は 0.878、最良の重みの baseline は 0.890 で、差は $-0.012$ である。正解式が 2 本以上のケースでは、差は $+0.10$〜$+0.18$ である。入力変数の数で分けると、差は 5 個以下で $+0.02$、6〜10 個で $+0.09$、11 個以上で $+0.16$〜$+0.19$ である。正解式と入力変数が多いケースほど、提案手法と baseline の差は大きい。

%======================================================================
\section{結果 2：向上の内訳と、モデルが頼っている特徴量}\label{sec:mech}
%======================================================================
ここでは、第 2 段の検索が、同じ重み $(0.7, 0.3)$ の baseline から $\rk$ を $+0.222$ 上げた内訳を調べる（設定 A）。以下の比較では、特徴量の組ごとに MLP を訓練し直し、同じ 10 通りの分割で評価した平均を示す。
\begin{itemize}[leftmargin=1.5em]
  \item \textbf{(i) 重みの学習}：第 1 段の点数の 2 項（文章類似度 $x_1$ と変数の一致度 $x_2$）だけを特徴量として MLP を訓練すると、$\rk$ は 0.521 から 0.601 に上がる（$+0.080$）。
  \item \textbf{(ii) 基本の特徴量}：2 個に残りの基本の特徴量 5 個（$x_3$〜$x_7$）を加えて 7 個にすると、0.601 から 0.693 に上がる（$+0.092$）。5 個のうち式の特化度 $x_6$ だけを加えても 0.691（$+0.090$、$p = 0.00195$）になり、この段階の向上はほとんど式の特化度による。
  \item \textbf{(iii) 集合の特徴量}：基本の 7 個に集合の 3 個（$x_8$〜$x_{10}$）を加えると、0.693 から 0.743 に上がる（$+0.050$、$p = 0.00195$）。
\end{itemize}
集合の 3 個を 1 個ずつ基本の 7 個に加えて訓練し直すと、補完性 $x_8$ で $+0.052$、一貫性 $x_9$ で $+0.050$（どちらも $p = 0.00195$）、分野ラベルの一致度（集合）$x_{10}$ で $-0.0001$（$p = 0.98$）であった。3 個すべてを加えたときの $+0.050$ は、補完性か一貫性の 1 個を加えたときと同程度であり、2 つの特徴量の効果は重なっている。

集合の特徴量による向上 $+0.050$ は、ケースによって異なる。正解式が 1 本のケースでは $+0.002$ であり（比べる相手の式がない）、3 本のケースで最大の $+0.102$ になる。入力変数の数で分けると、11〜15 個では $+0.078$ だが、16 個以上では $+0.035$ に下がる。この原因は \S\ref{sec:greedy} で調べる。

各特徴量への依存の度合いは、訓練し直さない方法でも調べた。訓練済みのモデル（10 個の特徴量）を固定したまま、テストのケースごとに、ある特徴量の値を候補どうしで無作為に入れ替え（シャッフル）、$\rk$ がどれだけ下がるかを測る（permutation feature importance, PFI）。入れ替えは 20 回行なって平均した。図~\ref{fig:pfi} にシャッフル後の $\rk$ を示す。変数の重複率を測る 6 特徴量（$x_2$、$x_4$〜$x_6$、$x_8$、$x_9$）をまとめてシャッフルすると、$\rk$ は 0.743 から 0.065 に下がる（低下量 0.678。低下量が 0 かを 10 通りの分割で調べた $t$ 検定で $p < 10^{-4}$）。1 個ずつでは、式の特化度をシャッフルしたときの低下が最も大きい（0.743 から 0.126）。話題の近さを測る 4 特徴量（$x_1$、$x_3$、$x_7$、$x_{10}$）は、まとめてシャッフルしても 0.737 であり、低下量 0.006 は 0 と有意に異ならない（$p = 0.10$）。特徴量を足して訓練し直す方法（アブレーション）と、訓練済みモデルの入力を入れ替える方法（PFI）は、同じ結論になった。効いているのは変数の重複率を測る特徴量であり、話題の近さを測る特徴量ではない。

\begin{figure}[t]\centering
\includegraphics[width=0.86\linewidth]{figure/fig_pfi_recall_ja.pdf}
\caption{特徴量の値をケース内で候補どうしシャッフルしたときの $\rk$（設定 A、提案手法、第 1 段の重み $(0.7, 0.3)$）。棒は 10 通りの分割の平均、誤差棒は $\rk$ の低下量の分割間の標準偏差、点線はシャッフル前の $\rk$（0.743）。各群の先頭の濃い棒は、その群の特徴量をまとめてシャッフルした結果。}
\label{fig:pfi}
\end{figure}

%======================================================================
\section{結果 3：逐次選択で、正解式の多いケースの向上が回復する}\label{sec:greedy}
%======================================================================
\S\ref{sec:mech} では、集合の特徴量による向上が、入力変数の多いケースで小さくなった。原因として、一括の並べ直しが参照集合を第 1 段の上位 5 件に固定しているため、正解式の多いケースでは、実際に上位に並ぶ式と参照集合が食い違うことを考えた。そこで、参照集合を、それまでに選んだ式に置き換える逐次選択（\S\ref{sec:seqmethod}）を使って、この原因だけを取り除いた。比較したのは、(1) 一括の並べ直し、(2) 逐次選択（推論のみ）、(3) 逐次選択（訓練も逐次）の 3 通りであり、同じ 10 通りの分割で訓練と評価を行なった。

この比較は、候補の件数を $N_1 = 200$ に広げて行なった。逐次選択が直すのは第 2 段の並べ方であるが、第 2 段は候補の中でしか並べ直せない。設定 A では、$N_1 = 50$ のとき正解式の 8.0\% が候補に入らないのに対し、$N_1 = 200$ では 2.2\% しか漏れないので、並べ方の違いが結果に現れやすい。この節の値は、表~\ref{tab:overall}（$N_1 = 50$）とは比べない。$N_1 = 50$ でも同じ比較を行なった（\ref{app:k50}）。

設定 B では、$\rk$ は (1) 0.724、(2) 0.756、(3) 0.765 の順になった。(2) は (1) を $+0.032$ 上回り（$p = 0.00195$）、(3) は (2) をさらに $+0.009$ 上回った（$p = 0.037$）。向上は、正解式の多いケースに集中した。正解式が 8 本以上のケースでは、(3) は (1) を $+0.056$（$p = 0.00195$）、(2) を $+0.013$（$p = 0.020$）上回った。設定 A でも (1) 0.757、(2) 0.787、(3) 0.797 と同じ順になったが、(3) と (2) の差 $+0.010$ は有意でなかった（$p = 0.13$）。設定 A には、参照集合が役に立ちにくい正解式 1〜2 本のケースが多いためと考えられる。

候補を $N_1 = 50$ に絞ったままでも、3 通りの順は同じであった（設定 B で (1) 0.689、(2) 0.709、(3) 0.717、設定 A で 0.743、0.762、0.765）。(3) は (1) を、設定 B で $+0.028$（$p = 0.00195$）、設定 A で $+0.022$（$p = 0.0039$）上回った。ただし、正解式が 8 本以上のケースでの (3) と (1) の差は、設定 B で $+0.021$、設定 A で $+0.013$ であり、$N_1 = 200$ の $+0.056$ と $+0.048$ より小さい。設定 A では、(3) と (2) の差も $+0.002$ で有意でなかった（$p = 0.49$）。$N_1 = 50$ では正解式の一部が候補に入らず、並べ方を変えても上位に来ないためと考えられる（\ref{app:k50}）。

%======================================================================
\section{結果 4：正解式の数を与えなくても選択を止められる}\label{sec:dof}
%======================================================================
自由度ゼロ停止（$K$ を与えない）を、同じ逐次選択（訓練も逐次、$N_1 = 200$）で上位 $K$ 本を採る場合（$K$ を与える）と比べた。完全一致率は、$K$ を与えた場合と同じであった（設定 A で 0.526、設定 B で 0.368。10 通りの分割の延べ 7{,}026 件のうち、判定が食い違ったケースは 0 件）。正解式をすべて選んだ時点で初めて停止条件を満たし、それより前には満たさなかったためである。集合 F1 は、$K$ を与えた場合よりわずかに低かった（設定 A で $-0.008$、設定 B で $-0.023$。対応のある $t$ 検定で $p = 0.006$ と $0.0001$）。

停止条件の充足率は、自由度ゼロ停止で 0.936（設定 A）と 0.883（設定 B）、$K$ を与えた場合で 0.865 と 0.781 であった。自由度ゼロ停止は条件を満たすまで選び続けるので、充足率が高いのは規則の定義による。条件を満たしても解けるとは限らないことは、\S\ref{sec:dofstop} で述べたとおりである。20 本選んでも条件を満たさなかったケースは 6.4\%（設定 A）と 11.7\%（設定 B）であり、すべて合成のケースで、その 98\%（設定 A）と 99\%（設定 B）では、正解式の一部が第 1 段の候補に入っていなかった。止まる位置の誤りは中央値 0 本で、誤るときはほとんどが多めに選ぶ方向であった（設定 B で誤った 545 件はすべて多め、設定 A では誤ったケースのうち 1 件だけが少なめ）。

注記：本データでは、正解式の数と出力変数の数がほとんど等しい（出典横断と合成では全ケースで等しい）ので、$K$ を当てること自体は難しくない。この結果が示すのは、$K$ を与えなくても選択を止められ、そのときの完全一致率が $K$ を与えた場合と同じであることであり、止まった式の集合が解けるモデルになることではない。

%======================================================================
\section{結果 5：話題の近さを測る特徴量はなぜ効かないか}\label{sec:topic}
%======================================================================
話題の近さを測る 4 特徴量は、どの形でも第 2 段の検索の精度を上げなかった。基本の 2 特徴量（$x_1$ と $x_2$）に分野ラベルの一致 $x_7$ を加えると $+0.002$（対応のある $t$ 検定で $p = 0.06$）、潜在意味類似度 $x_3$ を加えると $-0.002$（$p = 0.48$）、基本の 7 個に分野ラベルの一致度（集合）$x_{10}$ を加えると $-0.0001$（$p = 0.98$）であった。分野ラベルをまとめて粗くしても効かなかった。表記の違い（大文字と小文字、記号、複数形）をそろえて 1{,}958 種にした場合と、ラベルの TF-IDF ベクトルを k-means 法で 100 群と 40 群にまとめた場合に、$x_7$ を加えたときの向上は最大 $+0.0006$ で、すべて $p > 0.6$ であった。原因として、次の 3 つを確かめた。

\textbf{(1) 第 1 段の検索の後では、候補の話題がそろっている}：1 つの特徴量だけで、正解式と不正解式をどれだけ分けられるかを、AUC で調べた（MLP は使わない）。ケースごとに、正解式と不正解式の組をすべて作り、正解式の値の方が大きい組の割合（値が同じ組は 1/2 と数える）を求め、1{,}823 ケースで平均した。不正解式をデータベースの全式からとる場合と、第 1 段の候補 50 件の中からとる場合を比べる（表~\ref{tab:auc}）。潜在意味類似度の AUC は、全データベースに対しては 0.950 だが、候補 50 件の中では 0.594 に下がる。話題の近さで正解式を見分ける働きは、第 1 段の検索で使い切られている。

\begin{table}[t]\centering
\small
\caption{1 つの特徴量だけで正解式と不正解式を分けられる度合い（AUC。1{,}823 ケースの平均 $\pm$ 標準誤差）。列は、不正解式をとる範囲の違いを表す。式の特化度は全データベースに対して計算していない。}
\label{tab:auc}
\begin{tabular}{lcc}
\toprule
特徴量 & 全データベース $11{,}146$ 式 & 候補 50 件の中 \\
\midrule
変数の一致度 $x_2$ & $0.999 \pm 0.000$ & $0.947 \pm 0.002$ \\
式の特化度 $x_6$ & --- & $0.965 \pm 0.001$ \\
潜在意味類似度 $x_3$ & $0.950 \pm 0.003$ & $0.594 \pm 0.006$ \\
文章類似度 $x_1$ & $0.913 \pm 0.003$ & $0.437 \pm 0.007$ \\
\bottomrule
\end{tabular}
\end{table}

\textbf{(2) 候補の選び方が、見かけの逆相関を作る}：候補は、文章類似度と変数の一致度の重み付き和が大きい式だけを集めたものである。そのため候補の中では、文章類似度が高い式ほど、変数の一致度が低い式であることが多くなる（候補内の相関係数 $-0.404$。乱数の種 42 の分割、テストの 423 ケース）。正解式は変数の側で決まるので、候補の中では文章類似度がわずかに逆向きの手がかりになる（AUC 0.437）。

\textbf{(3) 分野ラベルの一致度（集合）は、ケースによって効く向きが逆で、打ち消し合っている}：PFI の値をケースごとに分けて見た。10 通りの分割のテストに現れるケースは延べ 4{,}594 件で、異なるケースは 1{,}592 件である（1 つのケースが 1〜8 回現れる）。同じケースを何回も数えると検定の $p$ 値が小さく出すぎるので、ケースごとに、現れた分割での値を平均して 1 ケース 1 つの値にしてから比べた。候補の中で $x_{10}$ の平均値が正解式の方で高いケース（1{,}169 件）では、シャッフルで $\rk$ が平均 0.012 下がる（特徴量が役立っている）。不正解式の方で高いケース（414 件）では、逆に平均 0.034 上がる（特徴量が精度を下げている）。2 群の差は Mann--Whitney の $U$ 検定で $p = 1.3 \times 10^{-8}$ である。2 群の値をケース数で重み付けて平均すると $(1{,}169 \times 0.0119 - 414 \times 0.0342)/1{,}583 = -0.0002$ となり、全体の PFI（$-0.0003$、図~\ref{fig:pfi}）とほとんど同じになる。モデルはこの特徴量をケースごとには使っているが、役立つケースと精度を下げるケースが打ち消し合って、全体では 0 に見える。分野ラベルをまとめても効かなかったので、この特徴量を生かすには、正解式の側で値が高くなるケースを増やす別の工夫が要る。

%======================================================================
\section{結果 6：誤りの分析と、第 1 段の重みの見直し}\label{sec:error}
%======================================================================
\subsection{誤りはどこで生まれるか}\label{sec:errwhere}
1 つの分割（乱数の種 42。テストのケース 423 件、正解式 1{,}579 本）で、提案手法（重み $(0.7, 0.3)$、$N_1 = 50$）の上位 $K$ 件に入らなかった正解式は 435 本あった（数は式の本数）。このうち、第 1 段の候補 50 件に入っていたのは 166 本（38.2\%）で、残る 269 本（61.8\%）は候補に入っていなかった。候補を 200 件にすると、上位 $K$ 件に入らなかった正解式は 373 本に減り、そのうち候補に入っていたのは 267 本（71.6\%）であった（候補に入らなかったのは 106 本）。$N_1 = 50$ では第 1 段の検索で、$N_1 = 200$ では第 2 段の並べ方で、多くの誤りが生まれている。

正解式をすべて上位 $K$ 件に並べられたケース（212 件）と、そうでないケース（211 件）を比べると、正解式の数は平均 2.13 本と 5.34 本、入力変数は 10.8 個と 17.8 個であり（どちらも Mann--Whitney の $U$ 検定で $p < 10^{-14}$）、後者は被覆率も 0.83 と低い（前者は 1.00）。候補に入ったのに上位 $K$ 件に入らなかった正解式と、それより上位にあった不正解式の組をすべて作ると、正解式の値の方が大きい組の割合は、潜在意味類似度で 0.516、変数の一致度で 0.375 であった。この誤りは、今ある特徴量では直せない。

\subsection{第 2 段の誤りの 2 つの型}\label{sec:errtypes}
\textbf{型 1：別の出典の同じ式}。標準的な式は複数の文献に載るので、正解式と同じ式を、別の出典から選ぶことがある（例：回分発酵のケースで、正解の $\mathrm{d}X/\mathrm{d}t = \mu X$ の代わりに、別の文献に載った同じ式が上位に来た）。評価では式を ID で区別するので、これは不正解と数えるが、利用者にとっては同じ式であり、使える式を提示できている。正解式 3{,}132 本のうち 386 本（12.3\%）は、空白などを除いて同じ文字列の式が別の出典にもある。本報告の検索手法の $\rk$ は、この型の分だけ低めに出ている。等価な式を正解と数える評価は今後の課題である。

\textbf{型 2：式の向き}。特徴量は変数の集合しか見ないので、同じ変数から別の変数を決める 2 本の式（例：線形化した反応器の $\mathrm{d}C_A'/\mathrm{d}t$ の式と $\mathrm{d}T'/\mathrm{d}t$ の式）は、すべての特徴量で同じ値になる。こちらは利用者にとっても誤りである。どの変数を決める式かを表す特徴量が要る。

\subsection{第 1 段の重みの見直し}\label{sec:redesign}
\S\ref{sec:errwhere} の分析から、第 1 段の検索で変数の一致度の重みを上げれば、被覆率（第 2 段の上限）が上がると予想した。そこで、第 1 段の重みを $(0.7, 0.3)$、$(0.5, 0.5)$、$(0.3, 0.7)$、$(0.0, 1.0)$ の 4 通りに変え、候補作り、訓練、評価のすべてを同じ重みでやり直した（設定 A、10 通りの分割）。

被覆率は予想どおり上がった（図~\ref{fig:redesign}(a)）。$N_1 = 50$ での被覆率は、重み $(0.7, 0.3)$ で 0.920、$(0.5, 0.5)$ で 0.961、$(0.3, 0.7)$ で 0.967、$(0.0, 1.0)$ で 0.961 であり、$(0.3, 0.7)$ はどの候補件数でも最も高い。文章類似度は、重みを下げて加えれば、変数の一致度だけのとき（0.961）より被覆率を上げる。

被覆率の向上は、ほとんどそのまま最終の精度の向上になった（図~\ref{fig:redesign}(b)）。提案手法の $\rk$ は、重み $(0.7, 0.3)$、$(0.5, 0.5)$、$(0.3, 0.7)$、$(0.0, 1.0)$ の順に 0.743、0.766、0.780、0.776 で、被覆率の順と同じ順になった。$(0.3, 0.7)$ は $(0.7, 0.3)$ を $+0.037$ 上回り、10 通りの分割すべてで上回った（$p = 0.00195$）。被覆率の増分 $+0.047$ のうち約 8 割が最終の精度に届いたことになる。向上は正解式が 8 本以上のケースに集中した（$+0.094$、$p = 0.00195$）。同じ重みの baseline も、0.521 から約 0.60 に上がる（0.602、0.600、0.603）。設定 B でも、重み $(0.3, 0.7)$ の提案手法は $(0.7, 0.3)$ を $+0.051$ 上回った（0.689 から 0.740、10 通りすべてで上回り $p = 0.00195$）。

候補を 200 件に広げた逐次選択（訓練も逐次）でも、重み $(0.3, 0.7)$ は $\rk$ を 0.797 から 0.810 に上げた（$+0.013$、$p = 0.0098$、10 通り中 9 通りで向上）。正解式が 8 本以上のケースでは $+0.035$（$p = 0.00195$）である。$N_1 = 200$ では重みによる被覆率の差が $+0.015$ まで縮むが、その約 9 割が最終の精度に届いた。第 2 段は、第 1 段が候補に入れた正解式をほとんど生かしている。残る誤りは \S\ref{sec:errtypes} の 2 つの型であり、候補の件数や重みでは直らない。

\begin{figure}[t]\centering
\includegraphics[width=0.98\linewidth]{figure/fig_stage1_redesign_ja.pdf}
\caption{第 1 段の重みを変える実験（設定 A）。(a) 重みごとの被覆率（1{,}823 ケースの平均 $\pm$ 標準誤差。点線は実験で使う $N_1 = 50$ と $200$）。(b) 候補作り、訓練、評価を同じ重みで行なったときの $\rk$（$N_1 = 50$。箱は 10 通りの分割の四分位、横線は中央値、点は各分割の値）。被覆率の順と最終の精度の順が一致し、文章 0.3、変数 0.7 が両方で最も高い。}
\label{fig:redesign}
\end{figure}

%======================================================================
\section{現時点の限界}\label{sec:limits}
%======================================================================
\begin{itemize}[leftmargin=1.5em]
  \item 合成ケースの説明文は LLM が書いたもので、人が書く依頼文と文体が違う可能性がある（式と入出力変数は規則で決まる）。単一出典の正解集合は、LLM が最小と判断した組である。
  \item 評価はデータベースの中で閉じており、データベースにない式は検索できない。
  \item LLM 直接生成との比較は各設定 50 件の 1 回だけの評価であり、等価の判定は LLM による。2 つの LLM の判定の一致は確かめたが、判定が正しいか（人手の判定との照合）と、同じ判定を繰り返したときのばらつきは調べていない。
  \item 検索手法の正解は ID で判定するので、別の出典の同じ式を選ぶと不正解と数える（\S\ref{sec:errtypes}）。
  \item 変数の重複率を測る 6 特徴量だけで、精度の大半が決まる（まとめてシャッフルすると $\rk$ が 0.678 下がる）。どのケースも入出力変数を正解式の変数記号から選んでおり、クエリと正解式が同じ表記を使うという本データの作り方も、これに効いている。表記の違うクエリに対して同じ精度が出るかは確かめていない。
  \item 逐次選択の訓練では正解式を参照集合にするが、推論では自分が選んだ式（誤りを含みうる）を参照集合にする。この食い違いは残っている。
  \item 自由度ゼロ停止は、要件 1〜4 のうち、要件 1 の出力変数の部分と、要件 2 を緩めた条件しか調べない。止まった式の集合が解けるとは限らない。
  \item 評価ケースのうち出典横断の 91 件（5.0\%）は、正解集合を与えても一部の未知の変数が決まらない（\S\ref{sec:problem}）。
  \item 検証用の分割を使っておらず、手法の定数は固定し、重みはテストの結果で比べた。\S\ref{sec:mech}〜\S\ref{sec:topic} の分析は、第 1 段の重み $(0.7, 0.3)$ で行なったものである。
\end{itemize}

%======================================================================
\section{まとめと今後の課題}\label{sec:summary}
%======================================================================
\textbf{まとめ}：提案手法は、第 1 段の検索で候補を絞り、第 2 段の検索で候補を 10 個の特徴量で並べ直す。最良の重みの古典的情報検索手法と比べ、$\rk$ を設定 A で 0.607 から 0.743（第 1 段の重みを $(0.3, 0.7)$ にすると 0.780）に、設定 B で 0.559 から 0.689（同 0.740）に上げた。向上は、正解式と入力変数の多いケースで大きい。効いているのは変数の重複率を測る特徴量（特に式の特化度）であり、話題の近さを測る特徴量は効かなかった。その原因は、第 1 段の検索の後で候補の話題がそろうこと、候補の選び方が作る見かけの逆相関、分野ラベルの一致度（集合）の効き方の打ち消し合いであった。逐次選択は、正解式が 8 本以上のケースで $\rk$ を $+0.056$ 上げた。自由度ゼロ停止は、$K$ を与えなくても、$K$ を与えた場合と同じ完全一致率で選択を止めたが、止まった式の集合が解けるとは限らない。誤りの過半が第 1 段の検索で生まれることを確かめ、第 1 段で変数の一致度の重みを上げたところ、予想どおり被覆率が上がり、$\rk$ が上がった。

\textbf{今後の課題}：
\begin{itemize}[leftmargin=1.5em]
  \item LLM 直接生成を全ケースで複数回評価する。等価の判定は、次の 2 点を確かめる。(1) 判定の正しさ：判定した組から無作為に選んだ組を人手で判定し、LLM の判定との一致率を求める。数式処理（SymPy など）で扱える組では、記号計算による判定とも比べる。(2) 判定のばらつき：同じ組を複数回判定させ、判定が変わる組の割合と、$\rk$ の回ごとの標準偏差を求める。
  \item 文献を与えた LLM との比較を用意する（全文献は LLM の入力の長さを超えるので、文献の与え方を検討する）。
  \item 第 1 段の重み $(0.3, 0.7)$ で、\S\ref{sec:mech}〜\S\ref{sec:topic} の分析をやり直す。重みなどの定数は、検証用の分割で選ぶ。
  \item 自由度ゼロ停止に要件 2〜4（等号の自由度、可解性、定数の一意性）を加え、止まった式の集合がどれだけ要件を満たすかを数える。
  \item 等価な式を正解と数える評価を用意し、どの変数を決める式かを表す特徴量を設計する。
\end{itemize}

\begin{thebibliography}{9}
\bibitem[Kato and Kano(2025)]{CombineEq}
S.~Kato and M.~Kano.
\newblock Automated physical model building from literature sources: Combining equations based on four predefined requirements.
\newblock \emph{Computers \& Chemical Engineering}, 200:109147, 2025.

\bibitem[Robertson and Zaragoza(2009)]{BM25}
S.~Robertson and H.~Zaragoza.
\newblock The probabilistic relevance framework: BM25 and beyond.
\newblock \emph{Foundations and Trends in Information Retrieval}, 3(4):333--389, 2009.

\bibitem[Seborg et~al.(2010)]{Seborg}
D.~E. Seborg, T.~F. Edgar, D.~A. Mellichamp, and F.~J. Doyle~III.
\newblock \emph{Process Dynamics and Control}.
\newblock John Wiley \& Sons, 3rd edition, 2010.

\bibitem[Wang et~al.(2022)]{E5}
L.~Wang, N.~Yang, X.~Huang, B.~Jiao, L.~Yang, D.~Jiang, R.~Majumder, and F.~Wei.
\newblock Text embeddings by weakly-supervised contrastive pre-training.
\newblock arXiv:2212.03533, 2022.
\end{thebibliography}

%======================================================================
\appendix
\def\thesection{付録 \Alph{section}}
%======================================================================
\section{比較手法の重みの比較}\label{app:weights}
表~\ref{tab:weights} に、古典的情報検索手法の重み $(w_{\mathrm{t}}, w_{\mathrm{v}})$ を 12 通りに変えた結果を示す。$\rk$ は設定 A、B とも $(0.4, 0.6)$ で最も高く、$w_{\mathrm{t}} \le 0.5$ の範囲ではほとんど変わらない（設定 A で 0.598〜0.607）。文章類似度だけ（$(1.0, 0.0)$）では 0.120（設定 A）と 0.153（設定 B）であり、説明文だけでは正解式を見つけにくい。

文章類似度を BM25（$k_1 = 1.5$、$b = 0.75$）や E5（\texttt{intfloat/e5-base-v2} による文の埋め込みベクトルのコサイン類似度）に置き換え、データベース全体で 0〜1 に正規化してから変数の一致度と混ぜた場合も、最良の重み $(0.1, 0.9)$ での $\rk$ は BM25 で 0.601（設定 A）と 0.555（設定 B）、E5 で 0.602 と 0.544 であり、TF-IDF の最良（0.607 と 0.559）と同程度である。提案手法（重み $(0.7, 0.3)$）は、どちらの設定でも、これらを 10 通りの分割すべてで上回った（$p = 0.00195$）。

\begin{table}[h]\centering
\small
\caption{古典的情報検索手法（TF-IDF の文章類似度と変数の一致度の重み付き和で全式を並べる）の重み別の結果（10 通りの分割の平均 $\pm$ 標準偏差）。太字は各設定で $\rk$ が最も高い重み。}
\label{tab:weights}
\footnotesize
\setlength{\tabcolsep}{2.5pt}
\begin{tabular}{@{}cccccccc@{}}
\toprule
 & \multicolumn{3}{c}{設定 A（1{,}823 件）} & & \multicolumn{3}{c}{設定 B（1{,}000 件）} \\
\cmidrule{2-4}\cmidrule{6-8}
重み $(w_{\mathrm{t}}, w_{\mathrm{v}})$ & $\rk$ & MAP & Recall@20 & & $\rk$ & MAP & Recall@20 \\
\midrule
$(1.0, 0.0)$ & $0.120 \pm 0.020$ & $0.148 \pm 0.025$ & $0.311 \pm 0.049$ & & $0.153 \pm 0.029$ & $0.170 \pm 0.030$ & $0.360 \pm 0.053$ \\
$(0.9, 0.1)$ & $0.215 \pm 0.030$ & $0.255 \pm 0.033$ & $0.503 \pm 0.065$ & & $0.254 \pm 0.031$ & $0.289 \pm 0.033$ & $0.516 \pm 0.048$ \\
$(0.8, 0.2)$ & $0.370 \pm 0.035$ & $0.433 \pm 0.045$ & $0.726 \pm 0.049$ & & $0.379 \pm 0.032$ & $0.423 \pm 0.034$ & $0.645 \pm 0.037$ \\
$(0.7, 0.3)$ & $0.521 \pm 0.045$ & $0.581 \pm 0.045$ & $0.836 \pm 0.023$ & & $0.466 \pm 0.037$ & $0.512 \pm 0.035$ & $0.721 \pm 0.032$ \\
$(0.6, 0.4)$ & $0.574 \pm 0.043$ & $0.652 \pm 0.044$ & $0.879 \pm 0.018$ & & $0.517 \pm 0.036$ & $0.574 \pm 0.036$ & $0.780 \pm 0.027$ \\
$(0.5, 0.5)$ & $0.602 \pm 0.042$ & $0.679 \pm 0.045$ & $0.898 \pm 0.017$ & & $0.546 \pm 0.039$ & $0.608 \pm 0.037$ & $0.814 \pm 0.027$ \\
$\mathbf{(0.4, 0.6)}$ & $\mathbf{0.607 \pm 0.043}$ & $0.692 \pm 0.046$ & $0.907 \pm 0.016$ & & $\mathbf{0.559 \pm 0.038}$ & $0.627 \pm 0.036$ & $0.827 \pm 0.027$ \\
$(0.3, 0.7)$ & $0.600 \pm 0.044$ & $0.691 \pm 0.046$ & $0.908 \pm 0.016$ & & $0.558 \pm 0.038$ & $0.630 \pm 0.037$ & $0.829 \pm 0.026$ \\
$(0.2, 0.8)$ & $0.600 \pm 0.044$ & $0.689 \pm 0.045$ & $0.906 \pm 0.016$ & & $0.554 \pm 0.038$ & $0.626 \pm 0.036$ & $0.826 \pm 0.023$ \\
$(0.1, 0.9)$ & $0.598 \pm 0.042$ & $0.685 \pm 0.044$ & $0.901 \pm 0.014$ & & $0.550 \pm 0.035$ & $0.620 \pm 0.035$ & $0.816 \pm 0.021$ \\
$(0.05, 0.95)$ & $0.599 \pm 0.041$ & $0.683 \pm 0.043$ & $0.899 \pm 0.013$ & & $0.550 \pm 0.032$ & $0.617 \pm 0.033$ & $0.812 \pm 0.022$ \\
$(0.0, 1.0)$ & $0.603 \pm 0.035$ & $0.679 \pm 0.039$ & $0.897 \pm 0.013$ & & $0.541 \pm 0.027$ & $0.605 \pm 0.029$ & $0.807 \pm 0.020$ \\
\midrule
BM25, $(0.1, 0.9)$ & $0.601 \pm 0.037$ & $0.692 \pm 0.041$ & $0.905 \pm 0.016$ & & $0.555 \pm 0.032$ & $0.624 \pm 0.035$ & $0.822 \pm 0.026$ \\
E5, $(0.1, 0.9)$ & $0.602 \pm 0.039$ & $0.684 \pm 0.039$ & $0.896 \pm 0.015$ & & $0.544 \pm 0.027$ & $0.611 \pm 0.031$ & $0.806 \pm 0.023$ \\
\bottomrule
\end{tabular}
\end{table}

\section{正解式の数と入力変数の数で分けた結果}\label{app:bucket}
表~\ref{tab:bucket} に、\S\ref{sec:overall} の差の内訳を示す（設定 A、10 通りの分割のテストのケースを合わせた延べ 4{,}594 件のケース平均。分割ごとのばらつきは示していない）。

\begin{table}[h]\centering
\small
\caption{正解式の数と入力変数の数で分けた $\rk$（設定 A）。提案手法は第 1 段の重み $(0.7, 0.3)$ の一括の並べ直し。差は提案手法から baseline を引いた値。}
\label{tab:bucket}
\begin{tabular}{@{}lrcccc@{}}
\toprule
 & 件数 & 提案手法 & baseline $(0.4, 0.6)$ & 差 & 差（baseline $(0.7, 0.3)$ との） \\
\midrule
正解式 1 本 & 418 & 0.878 & 0.890 & $-0.012$ & $+0.029$ \\
正解式 2 本 & 2{,}185 & 0.792 & 0.637 & $+0.154$ & $+0.227$ \\
正解式 3 本 & 318 & 0.690 & 0.528 & $+0.161$ & $+0.212$ \\
正解式 4 本 & 245 & 0.694 & 0.543 & $+0.151$ & $+0.250$ \\
正解式 5 本 & 242 & 0.745 & 0.563 & $+0.182$ & $+0.292$ \\
正解式 6 本 & 253 & 0.671 & 0.525 & $+0.146$ & $+0.252$ \\
正解式 7 本 & 212 & 0.616 & 0.483 & $+0.133$ & $+0.257$ \\
正解式 8 本 & 240 & 0.646 & 0.504 & $+0.142$ & $+0.262$ \\
正解式 9 本 & 244 & 0.643 & 0.504 & $+0.138$ & $+0.280$ \\
正解式 10 本 & 237 & 0.543 & 0.439 & $+0.104$ & $+0.262$ \\
\midrule
入力変数 5 個以下 & 485 & 0.850 & 0.828 & $+0.021$ & $+0.048$ \\
入力変数 6〜10 個 & 1{,}432 & 0.752 & 0.664 & $+0.088$ & $+0.173$ \\
入力変数 11〜15 個 & 916 & 0.746 & 0.553 & $+0.193$ & $+0.267$ \\
入力変数 16〜20 個 & 747 & 0.759 & 0.567 & $+0.192$ & $+0.289$ \\
入力変数 21 個以上 & 1{,}014 & 0.658 & 0.495 & $+0.163$ & $+0.283$ \\
\bottomrule
\end{tabular}
\end{table}

\section{逐次選択の比較（候補の件数 $N_1 = 50$ と $200$）}\label{app:k50}
表~\ref{tab:k50} に、\S\ref{sec:greedy} の 3 通りの比較を、候補の件数 $N_1 = 50$ と $200$ で示す（10 通りの分割の平均 $\pm$ 標準偏差）。$N_1 = 50$ の一括の並べ直しは、表~\ref{tab:overall} の提案手法（重み $(0.7, 0.3)$）と同じものである。

\begin{table}[h]\centering
\caption{逐次選択の 3 通りの $\rk$（第 1 段の重み $(0.7, 0.3)$）。差の括弧内は Wilcoxon の符号順位検定の $p$ 値。}
\label{tab:k50}
\footnotesize
\setlength{\tabcolsep}{3.5pt}
\begin{tabular}{@{}lcccc@{}}
\toprule
 & \multicolumn{2}{c}{設定 B} & \multicolumn{2}{c}{設定 A} \\
\cmidrule{2-3}\cmidrule{4-5}
 & $N_1 = 50$ & $N_1 = 200$ & $N_1 = 50$ & $N_1 = 200$ \\
\midrule
(1) 一括の並べ直し & $0.689 \pm 0.033$ & $0.724 \pm 0.041$ & $0.743 \pm 0.054$ & $0.757 \pm 0.056$ \\
(2) 逐次選択（推論のみ） & $0.709 \pm 0.030$ & $0.756 \pm 0.043$ & $0.762 \pm 0.049$ & $0.787 \pm 0.049$ \\
(3) 逐次選択（訓練も逐次） & $0.717 \pm 0.033$ & $0.765 \pm 0.043$ & $0.765 \pm 0.051$ & $0.797 \pm 0.053$ \\
\midrule
(2) $-$ (1) & $+0.020$ (0.0039) & $+0.032$ (0.00195) & $+0.019$ (0.0039) & $+0.030$ (0.00195) \\
(3) $-$ (2) & $+0.008$ (0.014) & $+0.009$ (0.037) & $+0.002$ (0.49) & $+0.010$ (0.13) \\
(3) $-$ (1) & $+0.028$ (0.00195) & $+0.041$ (0.00195) & $+0.022$ (0.0039) & $+0.040$ (0.00195) \\
(3) $-$ (1)、正解式 8 本以上 & $+0.021$ (0.00195) & $+0.056$ (0.00195) & $+0.013$ (0.00195) & $+0.048$ (0.00195) \\
\bottomrule
\end{tabular}
\end{table}

\end{document}
